\chapter{2012年重庆卷（理）}

\section{单选题}

\begin{problem}
在等差数列 \(\{a_{n}\}\) 中，\(a_{2} = 1\)，\(a_{4} = 5\)，则 \(\{a_{n}\}\) 的前 5 项和 \(S_{5} =\)（\quad）

\choices
    {\(7\)}
    {\(15\)}
    {\(20\)}
    {\(25\)}

\begin{answer}
B
\end{answer}

\begin{solution}
设等差数列的公差为 \(d\). 由 \(a_{4}-a_{2}=2d\)，得
\(2d=5-1=4\)，\(d=2\)
于是 \(a_{1}=a_{2}-d=1-2=-1\)，所以
\[
S_{5}=\frac{5}{2}\left(2a_{1}+4d\right)
=\frac{5}{2}\left(-2+8\right)=15
\]
故选 B.
\end{solution}
\end{problem}

\begin{problem}
不等式 \(\frac{x - 1}{2x + 1} \leqslant 0\) 的解集为（\quad）

\choices
    {\( \left(-\frac{1}{2},1\right] \)}
    {\( \left[-\frac{1}{2},1\right] \)}
    {\(\left(-\infty, -\frac{1}{2}\right) \cup [1, +\infty)\)}
    {\(\left(-\infty, -\frac{1}{2}\right] \cup [1, +\infty)\)}

\begin{answer}
A
\end{answer}

\begin{solution}
分式有意义的条件是 \(2x+1\ne0\)，即 \(x\ne-\frac12\). 令分子、分母分别为零，得到临界点 \(x=1\) 和 \(x=-\frac12\). 在三个区间上判断符号：
\[
\frac{x-1}{2x+1}>0\quad\text{当 }x<-\frac12\text{ 或 }x>1
\]
而在 \(-\frac12<x<1\) 时分式小于 \(0\). 由于 \(x=1\) 使分式等于 \(0\)，而 \(x=-\frac12\) 不在定义域内，解集为
\[
\left(-\frac12,1\right]
\]
故选 A.
\end{solution}
\end{problem}

\begin{problem}
对任意的实数 \(k\)，直线 \(y = kx + 1\) 与圆 \(x^{2} + y^{2} = 2\) 的位置关系一定是（\quad）

\choices
    {相离}
    {相切}
    {相交但直线不过圆心}
    {相交且直线过圆心}

\begin{answer}
C
\end{answer}

\begin{solution}
直线 \(y=kx+1\) 恒过定点 \((0,1)\). 该点到圆心 \(O(0,0)\) 的距离为 \(1\)，小于圆的半径 \(\sqrt2\)，所以定点在圆内，过该点的任意直线都与圆相交于两个点. 又直线不经过原点，因为将 \((0,0)\) 代入 \(y=kx+1\) 不成立. 因而直线与圆相交但不过圆心，故选 C.
\end{solution}
\end{problem}

\begin{problem}
\(\left(\sqrt{x} +\frac{1}{2\sqrt{x}}\right)^8\) 的展开式中常数项为（\quad）

\choices
    {\(\frac{35}{16}\)}
    {\(\frac{35}{8}\)}
    {\(\frac{35}{4}\)}
    {\(105\)}

\begin{answer}
B
\end{answer}

\begin{solution}
展开式的通项为
\(T_{r+1}=\mathrm C_{8}^{r}(\sqrt{x})^{8-r}\left(\frac{1}{2\sqrt{x}}\right)^{r} =\mathrm C_{8}^{r}2^{-r}x^{4-r}\)，\(r=0\)，\(1\)，\(\ldots\)，\(8\)
常数项要求 \(4-r=0\)，所以 \(r=4\). 因此常数项为
\[
\mathrm C_{8}^{4}2^{-4}=\frac{70}{16}=\frac{35}{8}
\]
故选 B.
\end{solution}
\end{problem}

\begin{problem}
设 \(\tan \alpha\)，\(\tan \beta\) 是方程 \(x^{2} - 3x + 2 = 0\) 的两根，则 \(\tan (\alpha + \beta)\) 的值为（\quad）

\choices
    {\(-3\)}
    {\(-1\)}
    {\(1\)}
    {\(3\)}

\begin{answer}
A
\end{answer}

\begin{solution}
由韦达定理，
\(\tan\alpha+\tan\beta=3\)，\(\tan\alpha\tan\beta=2\)
由于 \(1-\tan\alpha\tan\beta=-1\ne0\)，正切和公式可用，故
\[
\tan(\alpha+\beta)
=\frac{\tan\alpha+\tan\beta}{1-\tan\alpha\tan\beta}
=\frac{3}{1-2}=-3
\]
故选 A.
\end{solution}
\end{problem}

\begin{problem}
设 \(x\)，\(y \in \R\)，向量 \(\bs{a} = (x, 1)\)，\(\bs{b} = (1, y)\)，\(\bs{c} = (2, -4)\)，且 \(\bs{a} \perp \bs{c}\)，\(\bs{b} \parallel \bs{c}\)，则 \(|\bs{a} + \bs{b}| =\)（\quad）

\choices
    {\(\sqrt{5}\)}
    {\(\sqrt{10}\)}
    {\(2\sqrt{5}\)}
    {\(10\)}

\begin{answer}
B
\end{answer}

\begin{solution}
由 \(\bs a\perp\bs c\)，得
\(\bs a\cdot\bs c=2x-4=0\)，\(x=2\)
由 \(\bs b\parallel\bs c\)，得 \((1,y)=\lambda(2,-4)\). 由第一坐标 \(\lambda=\frac12\)，所以 \(y=-2\). 因而 \(\bs a=(2,1)\)，\(\bs b=(1,-2)\)，从而
\(\bs a+\bs b=(3,-1)\)，\(|\bs a+\bs b|=\sqrt{3^{2}+(-1)^{2}}=\sqrt{10}\)
故选 B.
\end{solution}
\end{problem}

\begin{problem}
已知 \(f(x)\) 是定义在 \(\R\) 上的偶函数，且以 2 为周期，则 \(f(x)\) 为 \([0,1]\) 上的增函数是“\(f(x)\) 为 \([3,4]\) 上的减函数”的（\quad）

\choices
    {既不充分也不必要的条件}
    {充分而不必要的条件}
    {必要而不充分的条件}
    {充要条件}

\begin{answer}
D
\end{answer}

\begin{solution}
设命题 \(P\) 为“\(f(x)\) 在 \([0,1]\) 上为增函数”，命题 \(Q\) 为“\(f(x)\) 在 \([3,4]\) 上为减函数”. 由周期性和偶性，对 \(x\in[3,4]\) 有
\[
f(x)=f(x-4)=f(4-x)
\]
其中 \(4-x\in[0,1]\).

若 \(P\) 成立，取 \(3\le x_{1}<x_{2}\le4\)，则 \(4-x_{1}>4-x_{2}\)，于是
\[
f(x_{1})=f(4-x_{1})\ge f(4-x_{2})=f(x_{2})
\]
所以 \(Q\) 成立. 反之，若 \(Q\) 成立，取 \(0\le u_{1}<u_{2}\le1\)，则 \(4-u_{2}<4-u_{1}\)，由 \(Q\) 得
\[
f(u_{2})=f(4-u_{2})\ge f(4-u_{1})=f(u_{1})
\]
所以 \(P\) 成立. 因此 \(P\) 是 \(Q\) 的充要条件，故选 D.
\end{solution}
\end{problem}

\begin{problem}
设函数 \( f(x) \) 在 \( \R \) 上可导，其导函数为 \( f'(x) \)，且函数 \( y = (1 - x)f'(x) \) 的图象如图所示，则下列结论中一定成立的是（\quad）

\bitmapfigure[max width=0.58\linewidth]{img/2012/chongqing_science/q8_fig1.png}

\choices
    {函数 \( f(x) \) 有极大值 \( f(2) \) 和极小值 \( f(1) \)}
    {函数 \( f(x) \) 有极大值 \( f(-2) \) 和极小值 \( f(1) \)}
    {函数 \( f(x) \) 有极大值 \( f(2) \) 和极小值 \( f(-2) \)}
    {函数 \( f(x) \) 有极大值 \( f(-2) \) 和极小值 \( f(2) \)}

\begin{answer}
D
\end{answer}

\begin{solution}
记 \(g(x)=(1-x)f'(x)\). 由图可知 \(g(x)\) 在 \(-2\)，\(1\)，\(2\) 处为零，并且符号依次为
\(g(x):\)，\(+\)，\(\ -\)，\(\ +\)，\(\ -\)
对应区间为 \((-\infty,-2)\)，\((-2,1)\)，\((1,2)\)，\((2,+\infty)\). 当 \(x<1\) 时 \(1-x>0\)，当 \(x>1\) 时 \(1-x<0\)，所以
\[
f'(x)>0\quad\text{在 }(-\infty,-2)\text{ 和 }(2,+\infty)
\]
\[
f'(x)<0\quad\text{在 }(-2,1)\text{ 和 }(1,2)
\]
因此 \(f(x)\) 在 \(-2\) 左增右减，在 \(2\) 左减右增，故有极大值 \(f(-2)\) 和极小值 \(f(2)\). 选项 D 正确.
\end{solution}
\end{problem}

\begin{problem}
设四面体的六条棱的长分别为 \(1\)，\(1\)，\(1\)，\(1\)，\(\sqrt{2}\) 和 \(a\)，且长为 \(a\) 的棱与长为 \(\sqrt{2}\) 的棱异面，则 \(a\) 的取值范围是（\quad）

\choices
    {\((0, \sqrt{2})\)}
    {\((0, \sqrt{3})\)}
    {\((1, \sqrt{2})\)}
    {\((1, \sqrt{3})\)}

\begin{answer}
A
\end{answer}

\begin{solution}
设长为 \(a\) 的棱为 \(AB\)，长为 \(\sqrt2\) 的异面棱为 \(CD\)，则其余四条棱 \(AC\)，\(AD\)，\(BC\)，\(BD\) 均长为 \(1\). 记 \(CD\) 的中点为 \(M\). 由 \(AC=AD=1\)，且 \(CM=DM=\frac{\sqrt2}{2}\)，得 \(AM\perp CD\)，并且

\[
AM=\sqrt{AC^{2}-CM^{2}}=\sqrt{1-\left(\frac{\sqrt2}{2}\right)^{2}}=\frac{\sqrt2}{2}
\]

因此，点 \(A\) 位于过 \(M\) 且垂直于 \(CD\) 的平面内，以 \(M\) 为圆心、半径 \(\frac{\sqrt2}{2}\) 的圆上. 同样，由 \(BC=BD=1\)，点 \(B\) 也位于这个圆上.
因此 \(AB\) 是此圆的一条弦，故
\[
0<a\le2\cdot\frac{\sqrt2}{2}=\sqrt2
\]
当 \(a=\sqrt2\) 时，\(A\)，\(B\) 为圆的两个对径点，四点 \(A\)，\(B\)，\(C\)，\(D\) 共面，不能构成四面体；当 \(a=0\) 时 \(A\) 与 \(B\) 重合，也不能构成四面体. 对于 \(0<a<\sqrt2\)，四点不共面，可以构成题设四面体. 所以
\[
0<a<\sqrt2
\]
故选 A.
\end{solution}
\end{problem}

\begin{problem}
设平面点集 \(A = \left\{(x, y) \mid (y - x)\left(y - \frac{1}{x}\right) \geqslant 0\right\}\)，\(B = \{(x, y) \mid (x - 1)^2 + (y - 1)^2 \leqslant 1\}\)，则 \(A \cap B\) 所表示的平面图形的面积为（\quad）

\choices
    {\(\frac{3}{4}\pi\)}
    {\(\frac{3}{5}\pi\)}
    {\(\frac{4}{7}\pi\)}
    {\(\frac{\pi}{2}\)}

\begin{answer}
D
\end{answer}

\begin{solution}
圆 \(B\) 的内部是圆心 \((1,1)\)、半径 \(1\) 的圆盘. 除去边界上的零测集，圆盘内有 \(x>0\)，\(y>0\)，于是条件可改写为
\[
(y-x)\left(y-\frac1x\right)\ge0
\quad\Longleftrightarrow\quad
(y-x)(xy-1)\ge0
\]
在交换 \(x\) 与 \(y\) 时，圆盘保持不变，而
\[
\left(y-x\right)(xy-1)
\quad\text{与}\quad
\left(x-y\right)(xy-1)
\]
符号相反. 因此，圆盘内满足严格不等式的一点与其关于直线 \(y=x\) 的对称点恰有一个属于 \(A\)，边界对面积没有影响. 所以 \(A\cap B\) 占圆盘面积的一半，面积为
\[
\frac12\cdot\pi\cdot1^{2}=\frac\pi2
\]
故选 D.
\end{solution}
\end{problem}

\section{填空题}

\begin{problem}
若 \( (1 + i)(2 + i) = a + bi \)，其中 \(a\)，\(b \in \R\)，\(i\) 为虚数单位，则 \( a + b = \)\fillinblank{}.

\begin{answer}
\(4\)
\end{answer}

\begin{solution}
\[
(1+i)(2+i)=2+i+2i+i^{2}=1+3i
\]
所以 \(a=1\)，\(b=3\)，从而 \(a+b=4\).
\end{solution}
\end{problem}

\begin{problem}
\(\displaystyle\lim_{n\to\infty}\frac{1}{\sqrt{n^2 + 5n} - n}=\)\fillinblank{}.

\begin{answer}
\(\frac{2}{5}\)
\end{answer}

\begin{solution}
有理化分母得
\[
\frac{1}{\sqrt{n^{2}+5n}-n}
=\frac{\sqrt{n^{2}+5n}+n}{5n}
=\frac{\sqrt{1+\frac5n}+1}{5}
\]
当 \(n\to\infty\) 时，\(\sqrt{1+\frac5n}\to1\)，故
\[
\displaystyle\lim_{n\to\infty}\frac{1}{\sqrt{n^{2}+5n}-n}
=\frac{1+1}{5}=\frac25
\]
\end{solution}
\end{problem}

\begin{problem}
设 \(\triangle ABC\) 的内角 \(A\)，\(B\)，\(C\) 所对的边分别为 \(a\)，\(b\)，\(c\)，且 \(\cos A = \frac{3}{5}\)，\(\cos B = \frac{5}{13}\)，\(b = 3\)，则 \(c =\fillinblank{}\).

\begin{answer}
\(\frac{14}{5}\)
\end{answer}

\begin{solution}
由于 \(A\)，\(B\) 是三角形内角，且余弦值为正，故
\(\sin A=\sqrt{1-\cos^{2}A}=\frac45\)，\(\sin B=\sqrt{1-\cos^{2}B}=\frac{12}{13}\)
又 \(C=\pi-(A+B)\)，所以
\[
\sin C=\sin(A+B)
=\sin A\cos B+\cos A\sin B
=\frac45\cdot\frac5{13}+\frac35\cdot\frac{12}{13}
=\frac{56}{65}
\]
由正弦定理，
\(\frac{c}{\sin C}=\frac{b}{\sin B}\)，\(c=3\cdot\frac{56}{65}\cdot\frac{13}{12}=\frac{14}{5}\)
\end{solution}
\end{problem}

\begin{problem}
过抛物线 \(y^{2}=2x\) 的焦点 \(F\) 作直线 \(l\) 交抛物线于 \(A\)，\(B\) 两点，若 \(|AB|=\frac{25}{12}\)，\(|AF|<|BF|\)，则 \(|AF|=\)\fillinblank{}.

\begin{answer}
\(\frac{5}{6}\)
\end{answer}

\begin{solution}
抛物线 \(y^{2}=2x\) 的焦点为 \(F\left(\frac12,0\right)\). 由于竖直弦 \(x=\frac12\) 的长度为 \(2\)，不等于题设的 \(\frac{25}{12}\)，故可设直线 \(l\) 为
\(y=k\left(x-\frac12\right)\)，\(k\ne0\)
与抛物线联立，得
\[
k^{2}\left(x-\frac12\right)^{2}=2x
\]
即
\[
k^{2}x^{2}-(k^{2}+2)x+\frac{k^{2}}4=0
\]
设交点横坐标为 \(x_{1}\)，\(x_{2}\)，则
\[
|x_{1}-x_{2}|=\frac{\sqrt{(k^{2}+2)^{2}-k^{4}}}{k^{2}}
=\frac{2\sqrt{k^{2}+1}}{k^{2}}
\]
因为两点在斜率为 \(k\) 的直线上，所以
\[
|AB|=|x_{1}-x_{2}|\sqrt{1+k^{2}}
=\frac{2(k^{2}+1)}{k^{2}}=\frac{25}{12}
\]
从而 \(k^{2}=24\). 代回交点方程，得
\(24x^{2}-26x+6=0\)，\((3x-1)(4x-3)=0\)
所以两个交点的横坐标为 \(\frac13\) 和 \(\frac34\). 由直线距离公式，
\(AF=\left|\frac13-\frac12\right|\sqrt{1+k^{2}}=\frac56\)，\(BF=\left|\frac34-\frac12\right|\sqrt{1+k^{2}}=\frac54\)
由 \(|AF|<|BF|\)，故 \(|AF|=\frac56\).
\end{solution}
\end{problem}

\begin{problem}
某艺校在一天的 \(6\) 节课中随机安排语文、数学、外语三门文化课和其他三门艺术课各 \(1\) 节，则在课表上的相邻两节文化课之间最多间隔 \(1\) 节艺术课的概率为\fillinblank{}.

\begin{answer}
\(\frac{3}{5}\)
\end{answer}

\begin{solution}
只考虑文化课和艺术课的类别排列即可. 三节文化课的位置从六个位置中选取，基本事件总数为
\[
\mathrm C_{6}^{3}=20
\]
设三节文化课之间的三段艺术课数分别为 \(g_{1}\)，\(g_{2}\)，并令两端的艺术课数为 \(g_{0}\)，\(g_{3}\). 则
\(g_{0}+g_{1}+g_{2}+g_{3}=3\)，\(g_{1}\)，\(g_{2}\in\{0,1\}\).
当 \((g_{1},g_{2})=(0,0)\) 时有 \(4\) 种；当恰有一个为 \(1\) 时有 \(3+3=6\) 种；当 \((g_{1},g_{2})=(1,1)\) 时有 \(2\) 种. 所以符合条件的排列有 \(4+6+2=12\) 种，所求概率为
\[
\frac{12}{20}=\frac35
\]
\end{solution}
\end{problem}

\section{解答题}

\begin{problem}
设 \( f(x)=a\ln x+\frac{1}{2x}+\frac{3}{2}x+1 \)，其中 \( a\in \R \)，曲线 \( y=f(x) \) 在点 \( (1,f(1)) \) 处的切线垂直于 \(y\) 轴.
\begin{enumerate}
\item 求 \(a\) 的值；
\item 求函数 \( f(x) \) 的极值.
\end{enumerate}
\begin{solution}
\begin{enumerate}
\item 函数的定义域为 \(x>0\)，且
\[
f'(x)=\frac{a}{x}-\frac{1}{2x^{2}}+\frac32
\]
曲线在点 \((1,f(1))\) 处的切线垂直于 \(y\) 轴，所以该切线为水平切线，故 \(f'(1)=0\). 因此
\(a-\frac12+\frac32=0\)，\(a=-1\)

\item 代入 \(a=-1\)，得
\[
f'(x)=-\frac1x-\frac{1}{2x^{2}}+\frac32
=\frac{3x^{2}-2x-1}{2x^{2}}
=\frac{(3x+1)(x-1)}{2x^{2}}
\]
在定义域 \(x>0\) 内，\(3x+1>0\)，\(2x^{2}>0\)，所以 \(f'(x)<0\) 当 \(0<x<1\)，\(f'(x)>0\) 当 \(x>1\). 于是 \(f(x)\) 在 \((0,1)\) 上递减，在 \((1,+\infty)\) 上递增，故 \(x=1\) 为极小值点. 极小值为
\[
f(1)=-\ln1+\frac12+\frac32+1=3
\]
函数没有最大值.
\end{enumerate}
\end{solution}
\end{problem}

\begin{problem}
甲，乙两人轮流投篮. 每人每次投一球. 约定甲先投且先投中者获胜，一直到有人获胜或每人都已投球 3 次时投篮结束. 设甲每次投篮投中的概率为 \(\frac{1}{3}\)，乙每次投篮投中的概率为 \(\frac{1}{2}\)，且各次投篮互不影响.
\begin{enumerate}
\item 求甲获胜的概率；
\item 求投篮结束时的甲投篮次数 \( \xi \) 的分布列与期望.
\end{enumerate}
\begin{solution}
\begin{enumerate}
\item 记甲、乙在第 \(k\) 次投篮投中的事件分别为 \(A_{k}\)，\(B_{k}\). 甲在第 \(k\) 次投中而获胜，要求前 \(k-1\) 轮双方都未投中，然后甲第 \(k\) 次投中. 因此
\(P(A_{k}\text{使甲获胜}) =\left(\frac23\cdot\frac12\right)^{k-1}\cdot\frac13 =\left(\frac13\right)^{k-1}\frac13\)，\(k=1\)，\(2\)，\(3\).
这些事件互斥，故甲获胜的概率为
\[
P(\text{甲获胜})=\frac13+\frac19+\frac1{27}=\frac{13}{27}
\]

\item \(\xi\) 的可能取值为 \(1\)，\(2\)，\(3\). 当第一轮结束时投篮结束的概率为
\[
P(\xi=1)=P(A_{1})+P(\overline{A_{1}}B_{1})
=\frac13+\frac23\cdot\frac12=\frac23
\]
第一轮双方均未投中后，第二轮结束的概率为
\[
P(\xi=2)=P(\overline{A_{1}}\,\overline{B_{1}})
\left[P(A_{2})+P(\overline{A_{2}}B_{2})\right]
=\frac23\cdot\frac12\cdot\left(\frac13+\frac23\cdot\frac12\right)
=\frac29
\]
若前两轮双方均未投中，则无论第三轮是否有人投中，投篮都在第三轮结束，所以
\[
P(\xi=3)=P(\overline{A_{1}}\,\overline{B_{1}}\,\overline{A_{2}}\,\overline{B_{2}})
=\left(\frac23\cdot\frac12\right)^{2}=\frac19
\]
因此分布列为
\begin{center}
\begin{tabular}{c|ccc}
\(\xi\) & \(1\) & \(2\) & \(3\) \\
\hline
\(P\) & \(\frac23\) & \(\frac29\) & \(\frac19\)
\end{tabular}
\end{center}
且
\[
E(\xi)=1\cdot\frac23+2\cdot\frac29+3\cdot\frac19=\frac{13}{9}
\]
\end{enumerate}
\end{solution}
\end{problem}

\begin{problem}
设 \( f(x)=4\cos\left(\omega x-\frac{\pi}{6}\right)\sin\omega x-\cos(2\omega x+\pi) \)，其中 \( \omega>0 \).
\begin{enumerate}
\item 求函数 \( y = f(x) \) 的值域；
\item 若 \( f(x) \) 在区间 \( \left[-\frac{3\pi}{2}, \frac{\pi}{2}\right] \) 上为增函数，求 \( \omega \) 的最大值.
\end{enumerate}
\begin{solution}
\begin{enumerate}
\item 令 \(\theta=\omega x\). 利用积化和差公式，
\[
4\cos\left(\theta-\frac\pi6\right)\sin\theta
=2\left[\sin\left(2\theta-\frac\pi6\right)+\sin\frac\pi6\right]
=2\sin\left(2\theta-\frac\pi6\right)+1
\]
又 \(\cos(2\theta+\pi)=-\cos2\theta\)，所以
\[
f(x)=1+2\sin\left(2\theta-\frac\pi6\right)+\cos2\theta
=1+\sqrt3\sin2\theta
=1+\sqrt3\sin(2\omega x)
\]
由于 \(-1\le\sin(2\omega x)\le1\)，函数值域为
\[
\left[1-\sqrt3,1+\sqrt3\right]
\]

\item 由上式
\[
f'(x)=2\sqrt3\,\omega\cos(2\omega x)
\]
因为 \(\omega>0\)，若 \(f(x)\) 在 \(\left[-\frac{3\pi}{2},\frac\pi2\right]\) 上为增函数，就必须在该区间内有 \(\cos(2\omega x)\ge0\). 当 \(x\) 从 \(-\frac{3\pi}{2}\) 变化到 \(0\) 时，\(2\omega x\) 从 \(-3\pi\omega\) 变化到 \(0\). 若 \(\omega>\frac16\)，则 \(-3\pi\omega<-\frac\pi2\)，这段自 \(-3\pi\omega\) 到 \(0\) 的区间必经过 \(\cos t<0\) 的部分，函数不可能在原区间上为增函数. 因而 \(\omega\le\frac16\).

当 \(\omega\le\frac16\) 时，对 \(x\in\left[-\frac{3\pi}{2},\frac\pi2\right]\)，有
\[
2\omega x\in\left[-3\pi\omega,\pi\omega\right]
\subseteq\left[-\frac\pi2,\frac\pi6\right]
\]
从而 \(\cos(2\omega x)\ge0\)，且函数在该区间上为增函数. 所以 \(\omega\) 的最大值为
\[
\frac16
\]
\end{enumerate}
\end{solution}
\end{problem}

\begin{problem}
如图，在直三棱柱 \(ABC - A_{1}B_{1}C_{1}\) 中，\(AB = 4\)，\(AC = BC = 3\)，\(D\) 为 \(AB\) 的中点.
\begin{enumerate}
\item 求点 C 到平面 \(A_{1}ABB_{1}\) 的距离；
\item 若 \(AB_{1} \perp A_{1}C\)，求二面角 \(A_{1} - CD - C_{1}\) 的平面角的余弦值.
\end{enumerate}

\bitmapfigure[max width=0.42\linewidth]{img/2012/chongqing_science/q19_fig1.png}
\begin{solution}
\begin{enumerate}
\item 在等腰三角形 \(ABC\) 中，\(D\) 为 \(AB\) 的中点，所以 \(CD\perp AB\)，且
\[
CD=\sqrt{AC^{2}-AD^{2}}=\sqrt{3^{2}-2^{2}}=\sqrt5
\]
直三棱柱的侧棱垂直于底面，故 \(AA_{1}\perp CD\)，\(BB_{1}\perp CD\). 因为 \(AB\) 与 \(AA_{1}\) 都在平面 \(A_{1}ABB_{1}\) 内，且 \(CD\) 同时垂直于这两条相交直线，所以 \(CD\perp\text{平面 }A_{1}ABB_{1}\). 因而点 \(C\) 到该平面的距离为 \(CD=\sqrt5\).

\item 取直角坐标系，使 \(D\) 为原点，\(DB\)，\(DC\)，\(DD_{1}\) 分别为 \(x\)，\(y\)，\(z\) 轴正方向，设棱柱的高为 \(h\). 则
\(A=(-2,0,0)\)，\(B=(2,0,0)\)，\(C=(0,\sqrt5,0)\)
\(A_{1}=(-2,0,h)\)，\(B_{1}=(2,0,h)\)，\(C_{1}=(0,\sqrt5,h)\)
由 \(AB_{1}\perp A_{1}C\)，得
\(\overrightarrow{AB_{1}}=(4,0,h)\)，\(\overrightarrow{A_{1}C}=(2,\sqrt5,-h)\)
\(\overrightarrow{AB_{1}}\cdot\overrightarrow{A_{1}C}=8-h^{2}=0\)，\(h=2\sqrt2\)
取 \(D_{1}=(0,0,h)\). 直线 \(DA_{1}\) 和 \(DD_{1}\) 都垂直于棱 \(CD\)，且分别在平面 \(A_{1}CD\) 和 \(C_{1}CD\) 内，所以二面角 \(A_{1}-CD-C_{1}\) 的平面角可取为 \(\angle A_{1}DD_{1}\). 因此
\(\overrightarrow{DA_{1}}=(-2,0,2\sqrt2)\)，\(\overrightarrow{DD_{1}}=(0,0,2\sqrt2)\)
\[
\cos\angle A_{1}DD_{1}
=\frac{\overrightarrow{DA_{1}}\cdot\overrightarrow{DD_{1}}}{|DA_{1}|\,|DD_{1}|}
=\frac{8}{(2\sqrt3)(2\sqrt2)}=\frac{\sqrt6}{3}
\]
\end{enumerate}
\end{solution}
\end{problem}

\begin{problem}
如图，设椭圆的中心为原点 \(O\)，长轴在 \(x\) 轴上，上顶点为 \(A\)，左，右焦点分别为 \(F_{1}\)，\(F_{2}\)，线段 \(OF_{1}\)，\(OF_{2}\) 的中点分别为 \(B_{1}\)，\(B_{2}\)，且 \(\triangle AB_{1}B_{2}\) 是面积为 4 的直角三角形.
\begin{enumerate}
\item 求该椭圆的离心率和标准方程；
\item 过 \(B_{1}\) 作直线 \(l\) 交椭圆于 \(P\)，\(Q\) 两点，使 \(PB_{2} \perp QB_{2}\)，求直线 \(l\) 的方程.
\end{enumerate}

\bitmapfigure[max width=0.42\linewidth]{img/2012/chongqing_science/q20_fig1.png}
\begin{solution}
\begin{enumerate}
\item 设椭圆方程为
\(\frac{x^{2}}{a^{2}}+\frac{y^{2}}{b^{2}}=1\)，\(a>b>0\)
焦点为 \(F_{1}(-c,0)\)，\(F_{2}(c,0)\)，其中 \(c^{2}=a^{2}-b^{2}\). 则
\(A=(0,b)\)，\(B_{1}=\left(-\frac c2,0\right)\)，\(B_{2}=\left(\frac c2,0\right)\)
三角形 \(AB_{1}B_{2}\) 为等腰三角形，直角只能在顶点 \(A\)，所以
\[
\overrightarrow{AB_{1}}\cdot\overrightarrow{AB_{2}}
=\left(-\frac c2,-b\right)\cdot\left(\frac c2,-b\right)
=b^{2}-\frac{c^{2}}4=0
\]
故 \(c^{2}=4b^{2}\)，从而 \(a^{2}=b^{2}+c^{2}=5b^{2}\). 又
\[
4=S_{\triangle AB_{1}B_{2}}=\frac12|B_{1}B_{2}|\cdot|OA|
=\frac12 cb=\frac12\cdot2b^{2}=b^{2}
\]
所以 \(b^{2}=4\)，\(c=4\)，\(a^{2}=20\). 因而离心率为
\[
e=\frac ca=\frac4{2\sqrt5}=\frac2{\sqrt5}
\]
椭圆的标准方程为
\[
\frac{x^{2}}{20}+\frac{y^{2}}4=1
\]

\item 由上问，\(B_{1}=(-2,0)\)，\(B_{2}=(2,0)\). 直线 \(l\) 不能为 \(x\) 轴，否则不满足垂直条件，故设
\(l:\)，\(x=my-2\)
将其代入椭圆方程，得
\(\frac{(my-2)^{2}}{20}+\frac{y^{2}}4=1\)，\((m^{2}+5)y^{2}-4my-16=0\)
设交点 \(P(x_{1},y_{1})\)，\(Q(x_{2},y_{2})\)，则 \(y_{1}\)，\(y_{2}\) 是上式的两根，所以
\(y_{1}+y_{2}=\frac{4m}{m^{2}+5}\)，\(y_{1}y_{2}=-\frac{16}{m^{2}+5}\)
由 \(x_{i}=my_{i}-2\)，有
\(\overrightarrow{B_{2}P}=(my_{1}-4,y_{1})\)，\(\overrightarrow{B_{2}Q}=(my_{2}-4,y_{2})\)
条件 \(PB_{2}\perp QB_{2}\) 等价于
\[
0=\overrightarrow{B_{2}P}\cdot\overrightarrow{B_{2}Q}
=(m^{2}+1)y_{1}y_{2}-4m(y_{1}+y_{2})+16
=\frac{16(4-m^{2})}{m^{2}+5}
\]
所以 \(m=\pm2\). 于是直线 \(l\) 为
\[
x-2y+2=0\quad\text{或}\quad x+2y+2=0
\]
\end{enumerate}
\end{solution}
\end{problem}

\begin{problem}
设数列 \(\{a_{n}\}\) 的前 \(n\) 项和 \(S_{n}\) 满足 \(S_{n + 1} = a_2S_n + a_1\)，其中 \(a_2\neq 0\).
\begin{enumerate}
\item 求证： \(\{a_{n}\}\) 是首项为 1 的等比数列；
\item 若 \(a_{2} > -1\)，求证：\(S_{n} \leqslant \frac{n}{2}(a_{1} + a_{n})\)，并给出等号成立的充要条件.
\end{enumerate}
\begin{solution}
\begin{enumerate}
\item 取 \(n=1\)，由题设得
\[
S_{2}=a_{2}S_{1}+a_{1}
\]
又 \(S_{1}=a_{1}\)，\(S_{2}=a_{1}+a_{2}\)，所以
\(a_{1}+a_{2}=a_{2}a_{1}+a_{1}\)，\(a_{2}(1-a_{1})=0\)
因 \(a_{2}\ne0\)，故 \(a_{1}=1\).

对 \(n\ge2\)，将
\(S_{n+1}=a_{2}S_{n}+a_{1}\)，\(S_{n}=a_{2}S_{n-1}+a_{1}\)
相减，得
\[
a_{n+1}=S_{n+1}-S_{n}=a_{2}(S_{n}-S_{n-1})=a_{2}a_{n}
\]
结合 \(a_{1}=1\)，可得
\[
a_{n}=a_{2}^{\,n-1}
\]
所以 \(\{a_{n}\}\) 是首项为 \(1\)、公比为 \(a_{2}\) 的等比数列.

\item 由上问，令 \(r=a_{2}\)，则 \(r>-1\)，\(r\ne0\)，且 \(a_{1}=1\)，\(a_{n}=r^{n-1}\).
\[
S_{n}=\sum_{k=0}^{n-1}r^{k}
\]
对 \(k=0\)，\(1\)，\(\ldots\)，\(n-1\)，有
\[
1+r^{n-1}-r^{k}-r^{n-1-k}
=\left(1-r^{k}\right)\left(1-r^{n-1-k}\right)\ge0
\]
这里，当 \(-1<r<0\) 时各因子均非负；当 \(0<r<1\) 时各因子均非负；当 \(r>1\) 时各因子均非正，乘积仍非负；\(r=1\) 时乘积为零. 对 \(k=0\) 到 \(n-1\) 求和，得到
\[
n\left(1+r^{n-1}\right)-2\sum_{k=0}^{n-1}r^{k}\ge0
\]
因此
\[
S_{n}\le\frac n2\left(1+r^{n-1}\right)
=\frac n2(a_{1}+a_{n})
\]

当 \(n=1\) 或 \(n=2\) 时，上式两端直接相等. 若 \(n\ge3\) 且 \(r\ne1\)，取 \(k=1\)，则 \(1-r\ne0\)，且 \(1-r^{n-2}\ne0\)，上面的求和严格为正，不能取等号. 当 \(r=1\) 时 \(a_{n}=1\)，显然对任意 \(n\) 都取等号. 所以等号成立的充要条件为
\[
n=1\quad\text{或}\quad n=2\quad\text{或}\quad a_{2}=1
\]
\end{enumerate}
\end{solution}
\end{problem}
